% ---------------------------------------------------------------------
%  Chapter 21 : The Theory-Builder UI
% ---------------------------------------------------------------------
\chapter{The Theory-Builder UI}\label{ch:ui}

Every other chapter of this manual assumes a \file{theories/<name>.theory.py}
file already exists --- a small Python module with a \code{build()} function that
chains a fluent builder (\code{TemporalTheoryBuilder(...)} or
\code{SpatialTheoryBuilder(...)}, then \code{.physical\_field(...)},
\code{.parameter(...)}, \code{.set\_action\_text(...)}, \code{.equation(...)},
\code{.build()}) into a model dict. The chapters on the theory specification
(Chapter~\ref{ch:theory-spec}) and on theory files (Chapter~\ref{ch:theory-files})
teach you to write those modules by hand. This chapter is about the alternative:
a point-and-click form, rendered \emph{inside the notebook}, that writes the file
for you.

The form is the class \code{TheoryUI}, and it lives in two files:
\file{pipeline/ui/main.py} (the form itself --- the ten-tab editor, the
validation sidebar, the buttons) and \file{pipeline/ui/widgets.py} (the small
reusable widget primitives the tabs are built from). The entire subsystem is
written on top of one external library, \term{ipywidgets}, the Jupyter project's
interactive-widget toolkit. If you have never used ipywidgets, this chapter
teaches you exactly as much of it as you need, from the ground up.

The single most important thing to understand before we start is what the UI
\emph{does not} do. \term{It does no physics.} It never enumerates a diagram,
never builds a propagator, never solves a saddle (with one small, optional
exception we will meet at the end). It is pure \emph{input specification}: it
collects the state of a form, assembles it into a plain Python dictionary --- the
\term{spec dict} --- and hands that dictionary to the serializer
\file{pipeline/theory\_serialize.py}, which renders it to a \code{.theory.py}
source file on disk. \textbf{Load} reverses the trip: it parses a saved file back
into a spec dict and repopulates the widgets. That is the whole job.

\begin{note}
Why a GUI at all, when the builder API is only a dozen lines? Because writing
those lines \emph{correctly} requires knowing the framework's conventions that
are easy to get wrong: the auto-generated response field \code{<name>t}, the
saddle parameter \code{<name>star}, the action grammar (\code{Dt} as a
time-derivative operator, \code{for i in pop} comprehensions, \code{Laplacian} in
spatial theories), and a pile of optional knobs (boundary conditions, Dyson
dressing order, operator-IR mode, stability analysis, run metadata). The form
surfaces every one of those as a labelled widget with help text, and a live
sidebar flags undeclared names and reserved-symbol collisions \emph{as you type}
--- so a typo is caught at authoring time instead of blowing up later with an
opaque \code{Sage} error deep in the engine.
\end{note}

\section{Where the UI sits in the pipeline}

The UI is the front door. A user opens the notebook
\file{notebooks/theory\_builder.ipynb}, runs one cell, fills in tabs, and clicks
\textbf{Save theory file}. From there the deliverable --- the \code{.theory.py}
file --- flows downstream exactly as if it had been hand-written. The data path
is short and one-directional on save, and reversible on load:

\begin{center}
\begin{tikzpicture}[
    node distance=6mm,
    box/.style={draw, rounded corners=2pt, align=center, font=\small,
                inner sep=5pt, fill=codebg},
    lbl/.style={font=\scriptsize\itshape, color=cmt},
    arr/.style={-{Stealth[length=2.4mm]}, thick}]
  \node[box] (nb) {\code{theory\_builder.ipynb}\\\code{TheoryUI().show()}};
  \node[box, below=of nb] (collect) {\code{TheoryUI.\_collect()}\\
        snapshot form $\rightarrow$ \textbf{spec dict}};
  \node[box, below=of collect] (ser) {\code{theory\_serialize}\\
        \code{save\_theory\_to\_file(spec, dir)}};
  \node[box, below=of ser] (file) {\file{theories/<slug>.theory.py}\\
        the deliverable on disk};
  \node[box, below=of file] (runner) {\code{theory\_runner.ipynb}\\
        \code{build()} $\rightarrow$ \code{compute\_cumulants(...)}};
  \draw[arr] (nb) -- (collect) node[lbl, midway, right] {user types in 10 tabs};
  \draw[arr] (collect) -- (ser) node[lbl, midway, right] {Save};
  \draw[arr] (ser) -- (file);
  \draw[arr] (file) -- (runner) node[lbl, midway, right] {the rest of the pipeline};
  % Load reverse arrow
  \draw[arr, dashed, color=Maroon] (file.west) to[out=180, in=180]
        node[lbl, midway, left, align=right] {Load:\\\code{load\_spec\_from\_file}\\
        $\rightarrow$ repopulate} (collect.west);
\end{tikzpicture}
\end{center}

\begin{itemize}
  \item \textbf{What feeds the UI:} the \emph{user} (keyboard input), and --- on
        \textbf{Load} --- an existing \file{theories/*.theory.py} file, parsed
        back into a spec dict by \code{load\_spec\_from\_file}.
  \item \textbf{What consumes the UI's output:} \code{pipeline.theory\_serialize}
        (the serializer that turns the spec dict into a \code{.theory.py} source
        string). The file in turn is consumed by the runner notebook (and the
        \code{nb\_support} engine), which calls \code{build()} to get the model
        dict and feeds it to \code{compute\_cumulants}. The UI also writes a
        one-line scratch file \file{.theories/.last\_built} so the
        \textbf{Open in runner notebook} button can hand the runner the
        just-saved name.
\end{itemize}

The boundary is clean: the UI's single output is the spec dict, and (via the
serializer) the \code{.theory.py} file. Everything downstream of that file is the
subject of the other chapters.

\section{Launching it}

The entire UI is two lines in a notebook cell. The package
\file{pipeline/ui/\_\_init\_\_.py} re-exports the class, so:

\begin{lstlisting}
from pipeline.ui import TheoryUI
TheoryUI().show()
\end{lstlisting}

That is the contents of the one cell in \file{notebooks/theory\_builder.ipynb}.
Constructing \code{TheoryUI()} builds every widget; \code{.show()} renders the
form into the cell's output area. The form then stays \emph{live}: typing in a box
or clicking a button runs Python on the kernel without re-executing the cell.

Two things happen the first time you call \code{show()}, and they are worth
knowing because they explain a small surprise later. Reading
\file{pipeline/ui/main.py:2311}:

\begin{lstlisting}
def show(self) -> None:
    global _CSS_INJECTED
    if not _CSS_INJECTED:
        display(HTML(_THEORY_BUILDER_CSS))
        _CSS_INJECTED = True
    try:
        self._refresh_validation()
    except Exception:
        pass
    display(self._root)
\end{lstlisting}

First, it injects an inline stylesheet --- one big \code{<style>} block in the
\code{tb-*} CSS namespace (\file{pipeline/ui/main.py:63}) --- but only \emph{once
per kernel session}, guarded by the module-level flag \code{\_CSS\_INJECTED}
(\file{pipeline/ui/main.py:244}). Second, it runs the validator a single time so
the sidebar shows real state on first render rather than placeholder text. Then it
\code{display}s the root container.

\begin{gotcha}
Because the CSS injection is guarded by a module-level flag, multiple
\code{TheoryUI()} instances in one kernel share the one stylesheet. If you edit
\code{\_THEORY\_BUILDER\_CSS} in the source, you must \emph{restart the kernel} to
see the change --- re-running the cell will not re-inject it. This is a deliberate
trade (the alternative is duplicate \code{<style>} blocks piling up on every
re-render), but it surprises people who are iterating on the styling.
\end{gotcha}

After you save, two attributes and one method form the public surface:
\code{ui.last\_saved} (the path of the last file written),
\code{ui.theories\_dir} (where saves land by default), and \code{ui.spec()}
(\file{pipeline/ui/main.py:2330}), which returns the current form state as a spec
dict without writing anything --- handy if you want to inspect or post-process the
dict in a later cell.

\section{ipywidgets in five minutes}

The rest of this chapter quotes ipywidgets code constantly, so here is the whole
mental model a physicist needs. Skip this section if you already know the library.

\term{ipywidgets} is a library of interactive HTML controls --- buttons, text
boxes, dropdowns, checkboxes, tabs, sliders --- that render inside a notebook
output cell and stay \emph{live}. The trick is that each widget is two halves kept
in sync: a Python object on the \emph{kernel} side and a DOM element on the
\emph{browser} side, connected over Jupyter's ``comm'' (communication) channel.
When you type in a box, the browser sends the new text to the kernel, the Python
object's \code{.value} attribute updates, and any callbacks you registered fire.
This is what lets a notebook cell host a full GUI with no web server of its own.

The conventional import alias is \code{W}, used in both source files
(\file{pipeline/ui/main.py:27}, \file{pipeline/ui/widgets.py:20}):

\begin{lstlisting}
import ipywidgets as W
from IPython.display import display, HTML
\end{lstlisting}

The handful of widget classes this code uses fall into four groups.

\paragraph{Layout containers.} \code{W.VBox([...])} and \code{W.HBox([...])} stack
their children vertically and horizontally. \code{W.Tab(children=[...])} is the
tab strip; \code{.set\_title(i, str)} names tab \code{i} and
\code{.selected\_index} reads or sets the active tab. \code{W.Layout(...)} is the
CSS-like layout object you attach to any widget --- \code{width='400px'},
\code{height='60px'}, or \code{display='none'} to hide it entirely.

\paragraph{Inputs.} \code{W.Text} (single line), \code{W.Textarea} (multi-line ---
used for the action and the seed box), \code{W.Dropdown} (single-select),
\code{W.Checkbox} (booleans), \code{W.ToggleButtons} (the Temporal/Spatial
switch), and the numeric trio \code{W.IntText} / \code{W.FloatText} /
\code{W.BoundedIntText} (where \emph{bounded} clamps to a \code{[min, max]}
range).

\paragraph{Output and display.} \code{W.HTML(value=...)} renders a static HTML
string --- every tab's help text, the readiness sidebar, and the cheat sheet are
\code{W.HTML} widgets. \code{W.Output(...)} is a \emph{capture region}: a panel
into which ordinary \code{print(...)} can be redirected. The UI's status panel is
one \code{W.Output}; the button handlers write to it with a
\code{with self.\_status:} context (\file{pipeline/ui/main.py:2642}), which
temporarily routes \code{stdout} into the panel, and \code{.clear\_output()} wipes
it. From \code{IPython.display}, \code{display(obj)} renders a rich object in the
current cell, \code{HTML(str)} wraps a raw HTML string, and \code{Javascript(str)}
ships a JS snippet to the browser (the \textbf{Open in runner} button uses it).

\paragraph{The event model.} Two registration calls appear everywhere:

\begin{lstlisting}
# fire a callback whenever a widget's .value changes:
self._w_theory_mode.observe(lambda _ch: self._on_mode_change(), names='value')
# fire a callback when a button is clicked:
self._btn_save.on_click(self._on_save)
\end{lstlisting}

\code{widget.observe(cb, names='value')} (\file{pipeline/ui/main.py:409}) calls
\code{cb} with a change-event object every time \code{.value} changes; the lambda
\code{lambda \_ch: ...} simply swallows that argument. \code{button.on\_click(cb)}
(\file{pipeline/ui/main.py:1376}) fires on click. Finally,
\code{widget.add\_class('tb-...')} attaches a CSS class to the widget's DOM node,
which is how the inline \code{tb-*} stylesheet finds its targets.

\begin{defn}
\term{observer} --- a callback registered with \code{widget.observe(cb,
names='value')} that fires whenever the widget's value changes. The UI registers
one observer, \code{\_mark\_changed}, on \emph{every} input, so any edit flips the
``unsaved changes'' flag and refreshes the validation sidebar. This is the
mechanism behind the live, type-as-you-go feedback.
\end{defn}

\section{The math behind each tab}

The UI does not implement the \msrjd{} field theory, but it is where the user
\emph{declares} it, so every tab maps onto a piece of the theory. A one-page
recap fixes what each tab is collecting and why the validator flags what it flags;
the full treatment is in Chapters~\ref{ch:theory-spec}--\ref{ch:enumeration}.

A stochastic equation of motion for a field $\phi$ (a scalar $x(t)$ for a temporal
theory, or a continuous field $\phi(x,t)$ for a spatial one) driven by noise
$\eta$,
\begin{equation}
  \Dt \phi \;=\; F[\phi] \;+\; \eta,
  \qquad \avg{\eta(t)\,\eta(t')} \;=\; 2D\,\delta(t-t'),
  \label{eq:ui-eom}
\end{equation}
is recast as a path integral by introducing a \term{response field} $\tilde\phi$
(Martin--Siggia--Rose / Janssen / De~Dominicis). The generating functional becomes
$Z = \int \mathcal{D}\phi\,\mathcal{D}\tilde\phi\,\ee^{-S[\phi,\tilde\phi]}$ with
the \msrjd{} \term{action}
\begin{equation}
  S[\phi,\tilde\phi] \;=\; \int \dd t\,\Big[\,\tilde\phi\,(\Dt\phi - F[\phi])
        \;-\; D\,\tilde\phi^{2}\,\Big].
  \label{eq:ui-action}
\end{equation}
This action is exactly the string the user types on the \textbf{Action} tab. The
placeholder text in the form (\file{pipeline/ui/main.py:1013}) reads:

\begin{lstlisting}[style=console]
phit * ((Dt + mu) * phi + eps * phi^3) - D * phit^2
\end{lstlisting}

which is $\tilde\phi\,(\Dt\phi + \mu\phi + \varepsilon\phi^{3}) - D\tilde\phi^{2}$
--- the action of $\dot\phi = -\mu\phi - \varepsilon\phi^{3} + \eta$ with
$\avg{\eta\eta} = 2D\delta$. Three syntactic pieces are visible:

\begin{itemize}
  \item \code{phit} --- the \term{response field} $\tilde\phi$. The UI never asks
        you to declare it: for every physical field \code{phi} you declare on the
        Fields tab, the framework auto-creates \code{phit} (response),
        \code{dphi} (fluctuation), and \code{phistar} (saddle). The action
        validator knows about all three (\file{pipeline/ui/main.py:1684}).
  \item \code{Dt} --- the \term{time-derivative operator} $\partial_t$, composed
        like an algebra element: \code{(Dt + mu)*phi} means
        $\partial_t\phi + \mu\phi$. It is one of the builtin identifiers the
        validator always accepts (\file{pipeline/ui/main.py:1653}).
  \item \code{- D phit\^{}2} --- the \term{Gaussian white-noise} term. Because the
        weight is $\ee^{-S}$, a term quadratic in $\tilde\phi$ corresponds to the
        \emph{second cumulant} of the noise: $-D\tilde\phi^{2}$ encodes
        $\kappa^{(2)} = 2D$ with a white ($\delta(\tau)$) kernel. A
        $-S_3\,\tilde\phi^{3}$ term would be a third white cumulant
        $\kappa^{(3)} = 3!\,S_3$.
\end{itemize}

With that, the role of each tab is:

\begin{description}
  \item[Model] the theory's name (which derives the filename), an optional
        description, and the Temporal/Spatial toggle.
  \item[Populations] a named group of \code{size N} identical units. Declaring a
        population \code{A} of size $N$ makes every field attached to it a
        length-$N$ vector and introduces a \code{for i in A} sum in the action.
        This is the ``$N$ coupled units'' case (oscillators, spins); a matrix
        parameter \code{w[i,j]} is the coupling between unit $i$ and unit $j$.
  \item[Fields] the physical fields $\phi$. The fluctuation split
        $\phi = \phi^{*} + \delta\phi$ expands the action around its saddle; you
        type only $\phi$ and the framework adds $\tilde\phi$, $\phi^{*}$,
        $\delta\phi$. In Spatial mode this tab also carries the
        \emph{spatial-structure} panel.
  \item[Parameters] the coefficients $\mu,\varepsilon,D,\dots$ appearing in
        $F[\phi]$ and the noise. Their \emph{shape} (scalar / vector indexed by
        one population / matrix indexed by two) follows the population structure.
  \item[Functions] non-polynomial transfer functions $f(\phi) = \tanh(a\phi)$ and
        the like. The framework Taylor-expands $f$ about the saddle of its first
        argument, $f(\phi) = f(\phi^{*}) + f'(\phi^{*})\delta\phi + \dots$, so you
        write \code{f[i](phi[i])} and the expansion is automatic.
  \item[Kernels] memory / convolution couplings $K * y = \int K(t-t')\,y(t')\,
        \dd t'$. The Fourier image $\tilde K(\omega)$ is preferred because the
        propagator builder works in frequency space.
  \item[Noise (CGF)] correlated, coloured, or cross-correlated noise declared as
        cumulant pieces $\avg{\eta(t)\eta(t')} = \text{coefficient}\times
        \text{kernel}(t-t')$. White $\to$ \code{dirac\_delta(tau)}; OU-coloured
        $\to$ \code{exp(-abs(tau)/tauc)}; cross-field $\to$ a coefficient like
        \code{rho*sqrt(D1*D2)} on two response legs.
  \item[Mean-field] the differential-algebraic-equation (DAE) residuals
        $\text{LHS} - \text{RHS} = 0$ solved at the saddle. \code{Dt} on the LHS
        marks a differential equation (set $\Dt\to 0$ at the saddle, kept symbolic
        for linear stability); no \code{Dt} means a purely algebraic equation. The
        solver finds \emph{all} roots and \code{fixed\_point\_index} picks which
        sorted root the diagrammatic expansion uses.
  \item[Defaults] the run metadata: \code{k} (number of external legs --- $k=2$ is
        the covariance / power spectrum), \code{max\_ell} (loop order --- \code{0}
        tree-level/LNA, \code{1} one-loop, \code{2} two-loop), the
        $(\code{tau\_max}, \code{tau\_step})$ lag grid, and the recommended
        external-field legs.
\end{description}

\begin{note}
\emph{Spatial structure} lives only on the Fields tab and only in Spatial mode.
Setting \code{spatial\_dim=1} makes $\phi$ a continuous field $\phi(x,t)$; the
action may then use the diffusion operator \code{Laplacian} multiplicatively
(\code{D*Laplacian*phi} $= D\Lap\phi$). \term{Derivative vertices} --- KPZ's
$(\partial_x h)^{2}$, Burgers, Model~B's $\Lap(\phi^{2})$ --- put a spatial
derivative \emph{inside} a nonlinearity and require the \term{Operator-IR} mode:
the action is then written with \code{Dt()} / \code{Lap()} / \code{Dx(phi, i)}
\emph{call} syntax instead of bare multiplication. We return to why this is
load-bearing in the Gotchas section.
\end{note}

None of this math is computed by the UI --- it is what the UI \emph{collects}. The
reader who wants the actual integrals should consult the theory-specification,
propagator, mean-field, and enumeration chapters.

\section{The widget primitives (\texttt{widgets.py})}

Before the form itself, look at the toolbox it is built from. The file
\file{pipeline/ui/widgets.py} (438 lines) defines a handful of small reusable
ipywidgets composites. Most of them are thin, but one --- \code{DynamicTable} ---
is the workhorse that every tabular tab is made of.

\subsection{Reading the system clipboard}

Several inputs carry a small clipboard-glyph button (rendered as a clipboard
emoji in the live form). Its reason for existing is a
genuine Jupyter wart, documented right in the source
(\file{pipeline/ui/widgets.py:23}): VS~Code's notebook renderer swallows
Ctrl/Cmd-V before it reaches an ipywidgets text input, so keyboard paste is broken
there. Button \emph{clicks} still work, and when the kernel runs on the same
machine as the editor (the normal case) Python can read the real system clipboard
--- so a clicked button reads it and fills the box directly.

\code{read\_system\_clipboard()} (\file{pipeline/ui/widgets.py:29}) returns the
local clipboard text or \code{None}. It tries \code{pyperclip} (an optional
cross-platform clipboard library) first, then falls back to OS-native command-line
tools picked by \code{sys.platform}: \code{pbpaste} on macOS;
\code{wl-paste}/\code{xclip}/\code{xsel} on Linux, in that preference order; and
PowerShell \code{Get-Clipboard} on Windows. For each candidate it skips the
command if \code{shutil.which} finds no binary, else runs it with
\code{subprocess.run(..., capture\_output=True, text=True, timeout=5)} and returns
\code{stdout} on a zero exit code.

\begin{gotcha}
The clipboard read happens on the machine the \emph{kernel} runs on, not your
browser. For a local Jupyter or VS~Code session that is your machine (correct);
for a \emph{remote} kernel it would read the server's clipboard (wrong, but
harmless --- you simply get the wrong text or a failure glyph). \code{paste\_button}
(\file{pipeline/ui/widgets.py:69}) handles the \code{None} case by briefly setting
its glyph to a cross mark and a tooltip suggesting \code{pip install pyperclip} or
a local kernel.
\end{gotcha}

\subsection{The scalar/vector/matrix and expression inputs}

Four small constructors round out the file. \code{vector\_input}
(\file{pipeline/ui/widgets.py:96}) builds a horizontal row of \code{FloatText}
widgets for a length-$N$ vector; \code{matrix\_input}
(\file{pipeline/ui/widgets.py:119}) an $n\times m$ grid; both expose a
monkey-patched \code{.get\_value()}. \code{expression\_input}
(\file{pipeline/ui/widgets.py:147}) is a single-line text box with a live
a live check-mark / cross-mark indicator driven by a \code{validator(text)}
callback.

\begin{gotcha}
\code{vector\_input}, \code{matrix\_input}, and \code{expression\_input} are
exported and importable, but the \emph{current form does not use them}. The form
takes vector and matrix defaults as a text cell on the Parameters tab, and does
its expression validation in the sidebar rather than per-field. They are kept as
reusable primitives; editing them will \emph{not} change the live UI. Of the bulk
import at \file{pipeline/ui/main.py:30}, only \code{textarea\_input},
\code{DynamicTable}, and \code{paste\_button} are actually called.
\end{gotcha}

The one composite the form \emph{does} use directly is \code{textarea\_input}
(\file{pipeline/ui/widgets.py:192}): a multi-line text area with a label and a
clipboard paste button, used for the action editor (\file{pipeline/ui/main.py:1008}).
It exposes the raw \code{Textarea} as \code{box.\_text\_w}, and the form reaches
into \code{\_text\_w} to attach observers and CSS classes.

\subsection{\texttt{DynamicTable} --- the spreadsheet primitive}

\code{DynamicTable} (\file{pipeline/ui/widgets.py:212}) is an add/remove-row
spreadsheet: each row is a dict of named values, with the columns declared up
front. Most of the form's tabs --- Populations, Fields, Parameters, Functions,
Kernels, Noise, Mean-field, and the external-fields table --- are a single
\code{DynamicTable}.

The constructor \code{\_\_init\_\_(self, columns, initial=None,
add\_label='+ add row')} (\file{pipeline/ui/widgets.py:245}) takes a list of
column-spec dicts. The recognized keys, from the class docstring
(\file{pipeline/ui/widgets.py:218}):

\begin{lstlisting}
{'name':             column key,
 'kind':             'text' | 'bool' | 'select' | 'int' | 'float',
 'options':          [...],               # static choices for 'select'
 'options_provider': callable() -> list,  # dynamic choices, re-queried on demand
 'default':          <value>,
 'width':            '120px',
 'placeholder':      'hint text',         # text columns
 'paste':            True}                # add a clipboard button beside a text cell
\end{lstlisting}

The factory \code{\_make\_widget} (\file{pipeline/ui/widgets.py:272}) dispatches on
\code{kind}: \code{text} $\to$ \code{W.Text}, \code{bool} $\to$ \code{W.Checkbox},
\code{select} $\to$ \code{W.Dropdown}, \code{int} $\to$ \code{W.IntText},
\code{float} $\to$ \code{W.FloatText}; an unknown kind raises \code{ValueError}.
Crucially, a \code{select} column's choices come from \code{options\_provider()} if
present, \emph{re-queried on demand}, else from the static \code{options} list ---
this is the hook that makes dropdowns track other tabs (see below). Every widget
gets an observer so edits propagate:

\begin{lstlisting}
w.observe(lambda _change: self._notify_change(), names='value')
\end{lstlisting}

\code{add\_row} (\file{pipeline/ui/widgets.py:309}) builds one widget per column;
for columns flagged \code{'paste': True} and \code{kind=='text'} it wraps the input
in an \code{HBox([w, paste\_button(w)])} --- but the \emph{value} widget is still
the bare \code{w}, so \code{get\_rows} reads it unchanged. Each row gets a
remove button (a cross-mark glyph).

\begin{gotcha}
The remove button is subtle. Its handler is built by
\code{\_make\_remove(len(self.\_row\_widgets))}, which \emph{captures} the row's
index at creation time --- but that captured index is effectively dead code. The
handler ignores it and instead scans \code{\_row\_boxes} for the box whose last
child \emph{is} this button (\file{pipeline/ui/widgets.py:332}), then pops that
index. This identity scan, not the stale captured index, is what makes removal
correct after earlier rows have already been deleted. If you ever refactor this,
keep the identity scan; the captured index is a trap.
\end{gotcha}

\code{get\_rows} (\file{pipeline/ui/widgets.py:365}) reads each row into a
\code{\{name: w.value\}} dict and is the table's single data-out path. It silently
\emph{drops rows whose \code{name} cell is blank}:

\begin{lstlisting}
def get_rows(self) -> list[dict]:
    out = []
    for w_dict in self._row_widgets:
        row = {name: w.value for name, w in w_dict.items()}
        if 'name' in row and not str(row['name']).strip():  # drop blank-name rows
            continue
        out.append(row)
    return out
\end{lstlisting}

\begin{gotcha}
Blank-row dropping is keyed on a column literally named \code{name}. The
external-fields table on the Defaults tab is keyed \code{field}, not \code{name},
so \code{get\_rows} does \emph{not} auto-drop its blank rows --- the metadata
collector re-filters on \code{field} instead (\file{pipeline/ui/main.py:2607}).
Any future table without a \code{name} column would silently keep its empty
starter rows unless it filters itself.
\end{gotcha}

Two more methods complete the picture. \code{set\_column\_visible(name, visible)}
(\file{pipeline/ui/widgets.py:375}) shows or hides one column --- header cell plus
every row's cell, current and future --- by toggling \code{layout.display} between
\code{''} and \code{'none'}; the form uses it to drop the \code{spatial\_dim}
column in Temporal mode. And \code{refresh\_dropdown\_options(col\_name)}
(\file{pipeline/ui/widgets.py:410}) re-queries a \code{select} column's
\code{options\_provider()} and pushes the new option list onto every existing
row's dropdown, preserving the prior selection when still valid:

\begin{lstlisting}
new_opts = list(provider())
for w_dict in self._row_widgets:
    w = w_dict.get(col_name)
    if not isinstance(w, W.Dropdown):
        continue
    prev = w.value
    try:
        w.options = new_opts
        if prev in new_opts:
            w.value = prev
        elif new_opts:
            w.value = new_opts[0]
    except Exception:
        # Some ipywidgets versions reject .options while .value is invalid:
        w.value = (new_opts[0] if new_opts else None)
        w.options = new_opts
\end{lstlisting}

This is how a population declared on the Populations tab instantly appears in the
index dropdowns on every other tab.

\begin{gotcha}
The \code{try/except} here is not defensive padding. Some ipywidgets versions raise
if you set \code{.options} while the current \code{.value} is momentarily not in
the new list; the fallback assigns \code{value} first, then \code{options}. If you
add a new population-aware dropdown anywhere, route its refresh through this method
--- never a raw \code{.options = ...} --- or you will hit that version-dependent
exception.
\end{gotcha}

Finally, \code{\_notify\_change} (\file{pipeline/ui/widgets.py:403}) calls every
registered \code{on\_change} callback inside a bare \code{except Exception: pass},
so one buggy observer cannot wedge the form. The trade is that a buggy validator
callback fails silently.

\section{The form (\texttt{main.py})}

\code{TheoryUI} (\file{pipeline/ui/main.py:289}) is the whole form. Its public
surface is small --- \code{\_\_init\_\_}, \code{show}, \code{spec}, \code{load},
the attributes \code{last\_saved} and \code{theories\_dir} --- and everything else
is private.

\subsection{Construction}

\code{\_\_init\_\_(self, theories\_dir=None)} (\file{pipeline/ui/main.py:296})
resolves the theories directory (default \code{\_DEFAULT\_THEORIES\_DIR}, which is
\code{<cwd>/../theories} --- since notebooks run from \code{notebooks/}, this
resolves to the repo's \file{theories/}), creates it if absent, and initializes
two pieces of bookkeeping state:

\begin{lstlisting}
self.last_saved: Optional[str] = None
self._dirty: bool = False          # unsaved-changes flag
self._loaded_extras: dict[str, dict[str, str]] = {
    'field_latex':    {},
    'function_latex': {},
    'kernel_latex':   {},
}
self._build_widgets()
\end{lstlisting}

\code{\_dirty} tracks unsaved edits. \code{\_loaded\_extras}
(\file{pipeline/ui/main.py:313}) is a per-kind \code{\{name: latex\}} carry-through
map: the \code{latex} column was removed from the UI on 2026-05-27, but older
theory files may still carry custom \code{latex} strings, so \code{load()} stashes
them here and \code{\_collect()} re-injects them. A load/edit/save round-trip thus
does not silently drop them.

\subsection{Population helpers --- the dropdown glue}

Three small methods make the cross-tab dropdowns work.
\code{\_pop\_names()} (\file{pipeline/ui/main.py:321}) returns the current non-blank
population names from the Populations table --- this is the \code{options\_provider}
that every population dropdown elsewhere is wired to. \code{\_pop\_size\_map()}
(\file{pipeline/ui/main.py:330}) returns \code{\{name: size\}} for populations with
a positive integer size. And \code{\_autofill\_default\_templates()}
(\file{pipeline/ui/main.py:345}) walks every Parameters row and, if its
\code{default} cell is empty and its index dropdowns name populations of known
size, fills a placeholder template (\code{'[, , ]'} for a length-3 vector,
\code{'[[, , ],[, , ]]'} for a $2\times3$ matrix) --- preserving any value the user
already typed. Note that this reaches directly into
\code{\_tbl\_parameters.\_row\_widgets} rather than \code{get\_rows()}, because it
needs to \emph{write} the widget \code{.value}.

\subsection{The ten tabs}

\code{\_build\_widgets()} (\file{pipeline/ui/main.py:381}) is the roughly
700-line constructor of the entire form. It builds, in order, the ten tab panels.
Each is either a few standalone widgets or a single \code{DynamicTable}:

\begin{longtable}{@{}>{\raggedright}p{0.16\textwidth}>{\raggedright\arraybackslash}p{0.78\textwidth}@{}}
\toprule
\textbf{Tab} & \textbf{What it builds} \\
\midrule
\endhead
1. Model & \code{\_w\_name} (\code{Text}), \code{\_w\_description}
  (\code{Textarea}), \code{\_w\_theory\_mode}
  (\code{ToggleButtons} of \code{['Temporal (ODE)', 'Spatial (PDE)']}); the
  toggle's observer calls \code{\_on\_mode\_change}. \\
2. Populations & \code{\_tbl\_populations} (name / size / description; no starter
  rows --- a scalar theory needs no populations). \\
3. Fields & \code{\_tbl\_physical} (name / population (select via
  \code{\_pop\_names}) / spatial\_dim / description; one starter row \code{phi}),
  plus the spatial-structure panel: \code{\_w\_spatial\_dim}
  (\code{BoundedIntText} 0--3), \code{\_w\_spatial\_apply} (a button that
  bulk-writes every row's \code{spatial\_dim}), \code{\_w\_boundary\_mode},
  \code{\_w\_boundary\_length}, \code{\_w\_initial\_mode}, \code{\_w\_operator\_ir}
  (\code{Checkbox}), \code{\_w\_dyson\_order}, \code{\_w\_reference\_diffusion}.
  These are grouped into \code{\_w\_spatial\_block}
  (\file{pipeline/ui/main.py:701}) so the toggle can hide them all at once. \\
4. Parameters & \code{\_tbl\_parameters} (name / index\_1 / index\_2 / domain /
  default (paste-enabled) / description; one starter row \code{mu}). The two index
  dropdowns are over \code{['\textemdash{}'] + populations} (the em-dash string is
  the ``no population'' sentinel); the default cell auto-templates when an index
  changes. \\
5. Functions & \code{\_tbl\_functions} (name / population / args / expression
  (paste) / description; no starter rows). \\
6. Kernels & \code{\_tbl\_kernels} (name / index\_1 / index\_2 / time\_expr
  (paste) / freq\_image (paste); no starter rows). \\
7. Noise & \code{\_tbl\_cgfs} (name / response\_legs / order / coefficient (paste)
  / kernel (paste); no starter rows). \\
8. Action & \code{\_w\_action = textarea\_input('Action S', ...)}
  (\file{pipeline/ui/main.py:1008}). \\
9. Mean-field & \code{\_tbl\_mfeqs} (lhs (paste) / rhs (paste) / population),
  \code{\_w\_fpi\_default} (\code{BoundedIntText} --- the default
  \code{fixed\_point\_index}), \code{\_w\_stability\_on} (\code{Checkbox}),
  \code{\_w\_seed\_box} (\code{Textarea} --- a multi-start Newton seed dict). \\
10. Defaults & \code{\_w\_k\_default}, \code{\_w\_ell\_default}
  (\code{BoundedIntText}), \code{\_w\_tau\_max}, \code{\_w\_tau\_step}
  (\code{FloatText}), and \code{\_tbl\_ext\_fields} (field / leaf\_index; two
  starter rows \code{phi[1], phi[2]}). \\
\bottomrule
\end{longtable}

The list of \code{(widget, title)} pairs is \code{\_all\_tabs}
(\file{pipeline/ui/main.py:1317}), and the \code{W.Tab} is composed from it with
numbered titles:

\begin{lstlisting}
self._all_tabs = [
    (tab_model, 'Model'), (tab_populations, 'Populations'),
    (tab_fields, 'Fields'), (tab_params, 'Parameters'),
    (tab_functions, 'Functions'), (tab_kernels, 'Kernels'),
    (tab_cgfs, 'Noise'), (tab_action, 'Action'),
    (tab_mfeqs, 'Mean-field'), (tab_defaults, 'Defaults'),
]
self._tabs = W.Tab(children=[w for w, _ in self._all_tabs])
for i, (_, t) in enumerate(self._all_tabs):
    self._tabs.set_title(i, f'{i + 1}. {t}')
\end{lstlisting}

After composing the tabs, \code{\_build\_widgets} wires the live plumbing
(\code{\_on\_populations\_changed} refreshes all dependent dropdowns; the parameter
table re-templates its defaults on an index change), builds the bottom bar,
attaches the readiness sidebar and CSS classes, registers \code{\_mark\_changed} on
every table and standalone widget, and finally calls \code{\_refresh\_save\_hint()}
and \code{\_on\_mode\_change(mark\_dirty=False)} for the initial render.

\begin{note}
The module docstring at \file{pipeline/ui/main.py:4} still calls this a
``9-section tab editor.'' That label is stale --- \code{\_all\_tabs} composes ten
tabs. When the prose and the code disagree, trust \code{\_all\_tabs}: there are
ten.
\end{note}

\section{Temporal vs.\ Spatial mode}

The Temporal/Spatial toggle on the Model tab does more than show and hide
controls; it actively rewrites form state so a saved theory carries only the
structure its kind allows. \code{\_is\_spatial\_mode()}
(\file{pipeline/ui/main.py:1522}) is simply \code{True} iff the toggle value starts
with \code{'Spatial'}. The work is in \code{\_on\_mode\_change(mark\_dirty=True)}
(\file{pipeline/ui/main.py:1525}):

\begin{itemize}
  \item \textbf{Switching to Temporal:} hide \code{\_w\_spatial\_block}, hide the
        \code{spatial\_dim} column, force every field's \code{spatial\_dim=0}, and
        \emph{clear} the operator-IR / Dyson / boundary / initial widgets, so a
        saved time-only theory carries no spatial structure.
  \item \textbf{Switching to Spatial:} show the block, show the column, bump
        \code{\_w\_spatial\_dim} to $\ge 1$, and \emph{clear the Kernels and Noise
        tables} --- spatial theories take only Gaussian white independent noise and
        have no convolution kernels, so neither can leak into a saved PDE theory.
\end{itemize}

Then \code{\_apply\_visible\_tabs()} (\file{pipeline/ui/main.py:1560}) rebuilds
\code{self.\_tabs.children} to the visible subset (Spatial hides the Kernels and
Noise tabs), renumbers the titles, and preserves the current selection when it
stays visible.

\begin{gotcha}
The clearing is \emph{intentional and silent, with no undo}. If you type kernels
or noise rows and then flip to Spatial, they are gone (\file{pipeline/ui/main.py:1548}).
If you set \code{spatial\_dim} and operator-IR and then flip to Temporal, those are
force-zeroed and cleared (\file{pipeline/ui/main.py:1534}). The rationale is that
each mode supports a strict subset of structure, but a user mid-edit can lose work.
Save before you flip if you are unsure.
\end{gotcha}

\section{The readiness sidebar}

The most valuable feature for a first-time author is the live validation sidebar.
It re-renders on every edit (via \code{\_mark\_changed}) and shows three things: a
per-tab readiness badge strip, a detailed list of findings, and a declared-names
cheat sheet. The renderer is \code{\_refresh\_validation}
(\file{pipeline/ui/main.py:2097}): a per-tab badge dot is red if any error, amber
if any warning, green otherwise --- note that \code{info}-severity findings do
\emph{not} trip a badge. The green ``No problems detected'' line only appears
once you have typed at least one field, parameter, or action, so a blank starter
form does not lie green.

\subsection{The action AST validator}

The flagship check is static validation of the action string. It is built on
Python's standard-library \code{ast} module, which parses Python \emph{source code}
into an \term{Abstract Syntax Tree} --- a tree of typed nodes (\code{ast.Name} = a
bare identifier, \code{ast.Call} = a function call, \code{ast.GeneratorExp} = a
\code{(... for i in ...)} comprehension) --- \emph{without executing it}. That last
point is the whole value: the action is analysed, never \code{eval}'d, so a typo is
caught at authoring time instead of detonating later inside Sage.

\code{\_validate\_action(self, action\_text, spec)}
(\file{pipeline/ui/main.py:1716}) first parses the string in expression mode and
catches a \code{SyntaxError} with its exact location:

\begin{lstlisting}
import ast
try:
    tree = ast.parse(text, mode='eval')
except SyntaxError as e:
    line = (e.lineno or 1); col = (e.offset or 0)
    return [('error', f'action syntax error at line {line}, col {col}: {e.msg}')]
\end{lstlisting}

It then walks the tree with a recursive \code{walk(node, bound)} that tracks which
names are \emph{bound} by comprehensions (\code{for i in A} binds \code{i}) and
lambda arguments, so an inner index is never falsely reported. Any \code{ast.Name}
that is neither a bound target nor a \emph{known} identifier is flagged as
undeclared (the first five are shown, then \code{+N more}); a \code{for i in X}
where \code{X} is not a declared population is flagged as a bad population; and
iterating over the legacy alias \code{pop} when no populations are declared is
called out as a foot-gun.

What counts as ``known'' is built by \code{\_action\_known\_names(spec)}
(\file{pipeline/ui/main.py:1663}). The set is the union of the builtins
(\code{\_ACTION\_BUILTINS}, \file{pipeline/ui/main.py:1651} --- operators
\code{Dt}, \code{Conv}; math \code{exp}, \code{log}, the trig family, \code{sqrt},
\code{abs}, \code{pi}, \code{e}, \code{I}, \code{oo}; \code{heaviside},
\code{dirac\_delta}, \code{sign}; indexing helpers \code{sum}, \code{range},
\code{len}), every declared parameter, every physical field \emph{with its three
auto companions}, every kernel and function name, and every population name (and
its \code{pop\_<name>} spelling, plus the legacy \code{pop}):

\begin{lstlisting}
for f in spec.get('physical_fields') or []:
    nm = (f.get('name') or '').strip()
    if not nm: continue
    names.add(nm)           # physical field itself
    names.add(f'{nm}t')     # response field (auto)
    names.add(f'd{nm}')     # fluctuation (auto)
    names.add(f'{nm}star')  # saddle
\end{lstlisting}

For spatial theories it additionally adds \code{Laplacian}; and \emph{only} when
the operator-IR checkbox is on, it adds the call forms \code{Lap}, \code{Dx},
\code{Gradient}, \code{GradX} (\file{pipeline/ui/main.py:1712}).

\begin{defn}
\term{AST (Abstract Syntax Tree)} --- the tree of Python syntax nodes produced by
\code{ast.parse}. Validating the action by walking this tree (rather than
executing the string) is what makes a typo like \code{hevyside} a tidy
``undeclared name'' finding in the sidebar instead of an opaque exception three
phases into the engine. Keeping \code{\_ACTION\_BUILTINS} explicit is what lets the
validator tell a real builtin from a misspelled one.
\end{defn}

\subsection{Cross-tab validation and the cheat sheet}

\code{\_validate(self)} (\file{pipeline/ui/main.py:1830}) walks the collected spec
and returns \code{\{tab\_index\_1based: [(severity, message), ...]\}} with severity
in \code{\{warn, error, info\}}. Its docstring is explicit about scope. It
\emph{does} check: undeclared-population references on fields, parameters, and
kernels (the ``silent \code{n\_populations=0} bug class''); the action AST
cross-reference; a Defaults-dict parse failure; reserved-name collisions; and
per-field mean-field-equation coverage for coupled theories. It deliberately does
\emph{not} nag about the theory name still being the default, stability being off,
or a blank action. The reserved names are \code{t}, \code{omega}, \code{Dt},
\code{delta\_D}, \code{delta\_Dp} everywhere, plus \code{k}, \code{Laplacian},
\code{x}, \code{y}, \code{z} in spatial theories. Critically, \code{\_validate}
wraps \code{\_collect()} in \code{try/except} so a malformed Defaults dict becomes
a \emph{finding} rather than a crash.

The cheat sheet, \code{\_build\_cheat\_sheet} (\file{pipeline/ui/main.py:2010}),
renders an HTML reference of all declared names by tab. Fields render as
\emph{scalars} (\code{n}, \code{nt}, \code{nstar}) when they have no valid
population, or \emph{indexed} (\code{n[i] for i in pop\_X}, response \code{nt[i]},
saddle \code{nstar[i]}) when they do --- showing you the exact spellings the action
will accept. Blank starter rows are filtered out.

\begin{gotcha}
Two blind spots are worth internalizing. First, the validator \emph{cannot} detect
a missing operator-IR toggle: it only allowlists the \code{Lap()}/\code{Dx()} call
names \emph{when the box is already on}, so a derivative-vertex action written
without ticking operator-IR validates cleanly yet computes the wrong physics (see
the next section). Second, the ``every field needs an equation'' check
(\file{pipeline/ui/main.py:1991}) is a word-boundary substring match on the LHS
text, not a parse, and only fires for $\ge 2$ fields (coupled theories) --- a field
whose name appears only inside a function call could pass spuriously.
\end{gotcha}

\section{From form to file: the spec dict}

Everything the form does on save funnels through one method. \code{\_collect(self)}
(\file{pipeline/ui/main.py:2335}) snapshots every tab into the \term{spec dict} ---
the flat dictionary the serializer consumes. It walks each table via
\code{get\_rows()}, drops blank rows, and maps the UI shapes onto the serializer's
shapes. The shape, assembled at \file{pipeline/ui/main.py:2498}, is:

\begin{lstlisting}
{
  'name':            str,                 # _w_name
  'populations':     [{'name','size','description'?}, ...],
  'n_populations':   int,                 # derived len(populations)
  'description':     str,
  'response_fields': [],                  # always empty -- auto-generated downstream
  'physical_fields': [{'name','population'?,'spatial_dim'?,'latex'?,'description'?}, ...],
  'parameters':      [{'name','indexed_by'?,'domain'?,'default'?,'description'?}, ...],
  'functions':       [{'name','args':[...],'expression','population'?,...}, ...],
  'kernels':         [{'name','indexed_by'?,'time_expr'?,'freq_image'?,'latex_name'?}, ...],
  'cgf_terms':       [{'name','response_field','response_legs'?,
                       'order','coefficient','kernel'}, ...],
  'action_text':     str,                 # the action textarea
  'equations':       [{'lhs','rhs','population'(or None)}, ...],
  'stability_analysis': bool,
  'default_fundamental': {},              # deliberately empty
  'metadata':        {...},               # see below
  # spatial-only keys, present ONLY when a field is spatial:
  'boundary':            {'mode': 'infinite'|'periodic', 'length'?: float|str},
  'initial':             {'mode': 'stationary'},
  'operator_ir':         True,            # only when ticked
  'dyson':               {'mode':'fixed','order':N},   # only when order>0
  'reference_diffusion': float,           # only when set
}
\end{lstlisting}

A small nested helper, \code{\_index\_list(row)}
(\file{pipeline/ui/main.py:2352}), translates the two index dropdowns (where
the em-dash string means none) into the \code{indexed\_by} list: \code{[]} for a scalar,
\code{[pop]} for a vector, \code{[pop1, pop2]} for a matrix. Parameter defaults are
parsed from their text cell into Python objects (with a raw-string fallback on
failure), the Noise rows are normalized by \code{\_normalize\_cgf\_rows}
(\file{pipeline/ui/main.py:249}), and the spatial blocks are emitted \emph{only
when at least one field is spatial} --- matching the serializer's
``emit-only-when-present'' rule.

\begin{gotcha}
Two spots use the \emph{raw} \code{eval} rather than \code{ast.literal\_eval} to
parse a default --- the \code{\_eval\_dict} helper inside \code{\_collect}
(\file{pipeline/ui/main.py:2342}) and the per-parameter default cell
(\file{pipeline/ui/main.py:2427}):
\begin{lstlisting}
obj = eval(text, {'__builtins__': {}}, {})
\end{lstlisting}
Passing \code{\{'\_\_builtins\_\_': \{\}\}} as globals strips access to
\code{open}, \code{\_\_import\_\_}, and friends, so a default like
\code{[[1,0],[0,1]]} evaluates fine but \code{\_\_import\_\_('os')} raises
\code{NameError}. This is \emph{weaker} than \code{ast.literal\_eval} --- it still
evaluates arbitrary operators --- but the threat model is the user's own typos on
the user's own machine, so it is acceptable. Do not mistake it for safe evaluation
of untrusted input. The Defaults seed-box, by contrast, \emph{does} use
\code{ast.literal\_eval} (\file{pipeline/ui/main.py:2586}).
\end{gotcha}

The run-metadata sub-dict comes from \code{\_collect\_metadata}
(\file{pipeline/ui/main.py:2570}):

\begin{lstlisting}
{
  'k_default': int, 'ell_default': int, 'tau_max': float, 'tau_step': float,
  'recommended_external_fields': [(field, leaf_index), ...],   # only if non-empty
  'fixed_point_index_default': int,                            # always written
  'seed_box_default': {var: (low, high), ...},                 # only if provided
}
\end{lstlisting}

Note that \code{fixed\_point\_index\_default} is written \emph{even when 0}
(\file{pipeline/ui/main.py:2620}), so a reload is unambiguous; most other optional
keys are emitted only when non-default. And \code{default\_fundamental} is
deliberately always emitted empty (\file{pipeline/ui/main.py:2544}) --- the
per-parameter \code{default=} declarations now carry the suggested values, and a
loaded file's \code{DEFAULT\_FUNDAMENTAL} is ignored on load.

\subsection{The serialized file}

When the form is filled like the placeholder --- one field \code{x} in population
\code{pop}; parameters \code{mu}, \code{eps}, \code{D}; the OU-quartic action ---
\code{\_collect()} yields a spec whose serialized form is exactly the on-disk file
the quickstart chapter loads:

\begin{lstlisting}
def build():
    return (
        TemporalTheoryBuilder('OU Quartic (white noise)')
        .population('pop', size=1)
        .physical_field('x', population='pop', description='variable')
        .parameter('mu', default=1.0, domain='positive')
        .parameter('eps', default=0.02, domain='positive')
        .parameter('D', default=1.0, domain='positive')
        .set_action_text('sum(xt[i]*((Dt+mu)*x[i] + eps*x[i]^3) - D*xt[i]^2 for i in pop)')
        .equation(lhs='(Dt+mu)*x[i]', rhs='-eps*x[i]^3', population='pop')
        .build()
    )
DEFAULT_FUNDAMENTAL = {'mu': 1.0, 'eps': 0.02, 'D': 1.0}
METADATA = {'k_default': 2, 'ell_default': 0,
            'recommended_external_fields': [('dx', 1), ('dx', 1)],
            'tau_max': 8.0, 'tau_step': 0.5, 'fixed_point_index_default': 0}
\end{lstlisting}

The serializer chooses \code{SpatialTheoryBuilder} versus
\code{TemporalTheoryBuilder} by whether any field carries \code{spatial\_dim} $\ge
1$. The saddle parameter \code{xstar} is \emph{not} emitted --- \code{build()}
auto-creates it from the field declaration. Compare this against the hand-written
file in Chapter~\ref{ch:quickstart}: they are identical, which is the point.

\begin{example}\label{ex:ui-roundtrip}
\textbf{The round trip.} Fill the form for the OU-quartic model and click
\textbf{Save}; you get the file above. Now in a fresh \code{TheoryUI()}, pick that
file in the Load dropdown and click \textbf{Load}: \code{load\_spec\_from\_file}
parses the file's \emph{source AST} (it never executes \code{build()}), the spec
dict is reconstructed, and \code{load(...)} clears and repopulates every tab. The
action text, equation RHS strings, function expressions, and kernel expressions
survive verbatim because they are stored as string literals; comments and
formatting are lost (it is an AST round-trip, not a textual one). Edit a parameter,
save again, and the file is regenerated --- closing the loop.
\end{example}

\section{Loading an existing theory}

\textbf{Load} reverses the save trip. \code{\_on\_load}
(\file{pipeline/ui/main.py:2769}) is gated by the dirty guard (below), refreshes
the dropdown so freshly-saved files appear, calls \code{load\_spec\_from\_file} to
parse the file into a spec dict, and hands it to \code{self.load(spec)}.

\code{load(self, spec\_or\_path)} (\file{pipeline/ui/main.py:2813}) is the big
repopulate-the-form routine, and most of its length is historical-wrinkle handling:

\begin{itemize}
  \item It strips the auto-prefix: a field stored as \code{dn} with
        \code{natural\_name='n'} displays as \code{n}.
  \item It stashes any \code{latex} strings into \code{\_loaded\_extras} for
        round-trip preservation.
  \item It \emph{auto-skips saddle parameters} (those with \code{mean\_field=True}
        or a name matching \code{<natural>star}) --- the framework re-creates them
        from the field declarations, so they must not be re-added as ordinary
        parameters.
  \item It maps \code{indexed\_by} back to the two index dropdowns, and converts a
        legacy \code{mf\_equations} shape (\code{\{saddle, rhs\}}) into the new
        \code{equations} shape (\code{\{lhs, rhs, population\}}) by stripping the
        trailing \code{star} and looking up the field's population.
  \item Finally it restores \code{stability\_analysis} and the spatial widgets,
        picks the Temporal/Spatial mode from the loaded fields and applies it, then
        drives every Defaults-tab widget from \code{metadata}.
\end{itemize}

\section{The bottom bar: Save, Preview, Pre-compute, and the rest}

\code{\_build\_bottom\_bar} (\file{pipeline/ui/main.py:1582}) composes the
save-path field, a derived-filename hint, the four primary action buttons, and the
destructive-action gate (two confirm checkboxes plus the Reset and Load rows). The
button handlers all write into the shared \code{\_status} panel.

\paragraph{Preview} \code{\_on\_preview} (\file{pipeline/ui/main.py:2640}) runs
\code{\_collect()} $\to$ \code{render\_theory\_file(spec)} and prints the resulting
\code{.theory.py} source into the status panel without writing anything --- a
look-before-you-leap.

\paragraph{Save} \code{\_on\_save} (\file{pipeline/ui/main.py:2650}) collects the
spec, resolves the target path (an explicit field value, or \code{theories\_dir},
joining a relative non-\code{.py} base), calls \code{save\_theory\_to\_file},
records \code{last\_saved}, clears \code{\_dirty} and both confirm checkboxes, and
prints the path plus the \code{THEORY\_NAME = <slug>} line you paste into the
runner. Both \code{\_collect()} and the write are wrapped so a failure prints a
humanized error rather than a traceback.

\paragraph{Reset} \code{\_on\_reset} (\file{pipeline/ui/main.py:2689}) is gated by
the dirty guard; on confirm it re-runs \code{\_build\_widgets()} to wipe state
cleanly, clears \code{\_dirty}, and re-\code{show()}s.

\paragraph{Open in runner} \code{\_on\_open\_in\_runner}
(\file{pipeline/ui/main.py:2231}) is the bridge to the runner notebook. If nothing
has been saved it prints a hint; otherwise it derives the slug from
\code{last\_saved}, writes it (one line) to \file{.theories/.last\_built}, and
\code{display}s a \code{Javascript} snippet:

\begin{lstlisting}
display(Javascript("window.open('theory_runner.ipynb', '_blank');"))
\end{lstlisting}

The runner reads \file{.theories/.last\_built} when its \code{THEORY\_NAME} is
blank, so the just-saved theory is pre-selected. It degrades to a printed
instruction when JS is unavailable (for example under an \code{nbconvert} export).

\subsection{The dirty guard}

Two pieces protect unsaved work. \code{\_mark\_changed}
(\file{pipeline/ui/main.py:2204}) sets \code{\_dirty=True} and refreshes the
sidebar; it is the observer registered on every input. \code{\_check\_dirty\_guard}
(\file{pipeline/ui/main.py:2217}) returns \code{True} (action allowed) only if the
form is clean, or if the matching confirm checkbox is ticked. Both Reset and Load
route through it, so they cannot silently destroy unsaved edits --- you must
explicitly tick ``discard unsaved changes'' or save first.

\subsection{Pre-compute: the one place the UI touches the engine}

\code{\_on\_precompute} (\file{pipeline/ui/main.py:2704}) is the single exception
to ``the UI does no physics.'' It lets you ask, before leaving the form, ``is this
theory structurally healthy?'' --- without writing a file. The sequence is:

\begin{lstlisting}
src = render_theory_file(spec)                 # render to source ...
ns: dict = {}
exec(compile(src, '<precompute-spec>', 'exec'), ns)  # ... exec in a fresh namespace ...
model = ns['build']()                          # ... call build() for an in-memory model
result = precompute(model, verbose=True)       # the structural sanity pass
\end{lstlisting}

\code{pipeline.precompute(model)} (a thin primitive in
\file{pipeline/\_precompute.py}) does one cheap structural pass: Taylor-expand the
action at order~2, verify the mean-field saddle, and build and cache the
propagator. It returns a result dict from which the handler prints a compact
verdict --- the mean-field check (a \code{PASS} mark when \code{mf\_check == 'PASS'}),
the solved saddle values, whether the propagator was \code{cached} or
\code{FAILED}, the cache directory, and the wall time. A green PASS means the
theory builds, has a real saddle, and has a constructible propagator --- the most
common authoring mistakes (an undeclared symbol, a saddle that does not solve, an
ill-posed bilinear part) are caught here, in seconds, before you ever open the
runner.

\begin{gotcha}
\textbf{Fork safety is not a concern here.} The macOS fork-crash documented in this
project's memory lives in the spatial and temporal \emph{integrator} paths, not the
UI. The UI runs entirely on the kernel's main thread with no multiprocessing, and
\code{Pre-compute} runs \code{precompute} in-process. You can click it freely in a
notebook on any platform.
\end{gotcha}

\section{Gotchas to keep in mind}

Most caveats have appeared in context above; this section collects the ones most
likely to bite, in rough order of severity.

\begin{itemize}
  \item \textbf{Operator-IR is load-bearing and silent.} Per the serializer
        (\file{pipeline/theory\_serialize.py:455}), without \code{.operator\_ir()}
        a re-saved KPZ/Burgers/Model-B theory \emph{silently loses} the
        \code{Dt()}/\code{Lap()}/\code{Dx()} lowering and computes the wrong
        physics. The UI emits it only when the checkbox is ticked, and the
        validator does \emph{not} detect a missing tick. If your action has a
        spatial derivative \emph{inside} a nonlinearity, you must tick Operator-IR.
  \item \textbf{Mode-switch data loss.} Temporal~$\to$~Spatial clears the Kernels
        and Noise tables; Spatial~$\to$~Temporal force-zeros every
        \code{spatial\_dim} and clears operator-IR / Dyson / boundary. Both are
        intentional and have no undo. Save before flipping.
  \item \textbf{The \code{eval} sandbox is weak.} Parameter defaults are parsed
        with \code{eval(text, \{'\_\_builtins\_\_': \{\}\}, \{\})}, which blocks
        name lookups but still evaluates operators. Fine for your own typos; not
        safe-evaluation of untrusted input.
  \item \textbf{\code{latex} carry-through is keyed by name.} A loaded \code{latex}
        string survives a round trip via \code{\_loaded\_extras}, but if you
        \emph{rename} the field it was attached to, the latex is dropped.
  \item \textbf{Clipboard reads the kernel's machine.} The clipboard buttons exist
        because VS~Code swallows Ctrl/Cmd-V; on a remote kernel they read the
        server's clipboard (or show a cross mark).
  \item \textbf{Open-in-runner needs JS and relative paths.}
        \code{window.open('theory\_runner.ipynb', ...)} assumes the runner sits next
        to the builder in the Jupyter file tree and that JS is allowed; otherwise it
        prints an instruction.
  \item \textbf{The exported \code{vector\_input} / \code{matrix\_input} /
        \code{expression\_input} are unused by the live form.} They are reusable
        primitives; editing them changes nothing you see in the UI.
\end{itemize}

\section{What you just learned, and where to go next}

The Theory-Builder UI is the only chapter in this manual that produces no numbers.
Its job is upstream of everything else: it is a typed, validated, point-and-click
authoring surface whose single output is a \code{.theory.py} file identical to one
you could write by hand. You saw the kernel\,$\leftrightarrow$\,browser widget
model that makes a live form possible inside a notebook cell; the
\code{DynamicTable} primitive that every tabular tab is made of; the ten tabs and
the \msrjd{} piece each one declares; the Temporal/Spatial toggle and what it
rewrites; the AST-based action validator that catches typos before they reach the
engine; the \code{\_collect()} method and the spec-dict shape it produces; the
serializer round trip via \code{load\_spec\_from\_file}; and the bottom-bar buttons,
including the in-process \code{Pre-compute} sanity pass.

From here the natural next step is to take a file the UI produced and run it.
Chapter~\ref{ch:quickstart} drives exactly such a file end to end; the theory-files
chapter (Chapter~\ref{ch:theory-files}) documents the \code{build()} /
\code{DEFAULT\_FUNDAMENTAL} / \code{METADATA} shapes the UI emits, from the
serializer's side; and the serializer's own internals --- \code{render\_theory\_file},
\code{save\_theory\_to\_file}, \code{load\_spec\_from\_file} --- live in
\file{pipeline/theory\_serialize.py}, the module on the far side of every Save and
Load button you just met.
